% Prose rendering of PrimeTensor/Analysis/Derivative.lean.
% Fragment: no \documentclass, no preamble, no document environment.

\section{Multiplicative differential germs}
\label{Analysis:Derivative:sec}

\leanfile{PrimeTensor/Analysis/Derivative.lean}

\subsection*{What this file establishes}

This is where the development starts doing calculus, and the module header is
emphatic about how: the multiplicative derivative is \emph{not} obtained by
translating the additive difference quotient.  Four notions are introduced.

\begin{itemize}[leftmargin=1.4em]
\item A \emph{perturbation} is a sequence of completed magnitudes tending to
  the pivot $1$.  It replaces ``$h \to 0$''.
\item The \emph{response} of $f$ at $x$ along $h$ is the sequence of output
  ratios $f(x \cdot h_n)/f(x)$.  It replaces the difference $f(x+h) - f(x)$
  --- and note that nothing is divided by the perturbation.
\item A \emph{tangent map} is a morphism of the multiplicative carrier.  It
  replaces the linear map.
\item \emph{Negligibility} of one germ relative to another is stated by
  demanding that the first reach a scale level an arbitrary prescribed number
  of refinements deeper than the second.  It replaces little-o.
\end{itemize}

Assembling these, $D$ is a derivative of $f$ at $x$ when it models every
pivot-approaching perturbation to first order.  The file closes by computing
the two derivatives that can be computed with no analysis at all: the identity
function has the identity tangent map, and every constant function has the
trivial one.

Design note~\ref{Analysis:Derivative:note:log} works out what all four notions
say in the logarithmic chart, which is the quickest way to see that they are
the familiar ones and not merely analogous to them.

\subsection{Advancing a scale}

\begin{definition}[\lean{Depth.advance}]\label{Analysis:Derivative:def:advance}
$\name{advance} : \name{Depth} \to \name{Depth} \to \name{Depth}$ is defined by
recursion on its second argument:
\[
  \name{advance}(\ell, \ctor{one}) = \ctor{succ}\,\ell,
  \qquad
  \name{advance}(\ell, \ctor{succ}\,g) = \ctor{succ}\,\name{advance}(\ell, g).
\]
\end{definition}

\begin{leancode}
def advance : Depth → Depth → Depth
  | level, .one => .succ level
  | level, .succ gain => .succ (advance level gain)
\end{leancode}

So $\size{\name{advance}(\ell,g)} = \size{\ell} + \size{g}$: the level is
refined $\size{g}$ times.  As the docstring says, there is deliberately no
advance by zero --- $\ctor{one}$ already means one refinement.  A consequence
worth recording now is that the vocabulary can express ``strictly finer by a
prescribed amount'' and cannot express ``no finer'': there is no way to write
the multiplicative analogue of a big-$O$ relation in this file, only of a
little-o one.

\subsection{Ratios in the completion}

\begin{definition}[\lean{MulReal.ratio}]\label{Analysis:Derivative:def:ratio}
$\name{ratio}(a,b) \defeq a \cdot b^{-1}$, the completed counterpart of the
ratio of \texttt{PrimeTensor/Order.lean}.
\end{definition}

\begin{lemma}[\lean{MulReal.ratio\_self}, \lean{inv\_one} and
  \lean{ratio\_mul\_base}]\label{Analysis:Derivative:lem:ratio}
$\name{ratio}(a,a) = 1$; $1^{-1} = 1$; and for all $x, h$,
\[
  \name{ratio}(x \cdot h,\; x) = h .
\]
\end{lemma}

\begin{leancode}
theorem ratio_mul_base (x h : MulReal) :
    ratio (x * h) x = h := by
  unfold ratio
  calc
    (x * h) * x⁻¹
        = h * (x * x⁻¹) := by
          rw [mul_comm x h]
          exact mul_assoc h x x⁻¹
    _ = h * 1 := by rw [mul_inv x]
    _ = h := mul_one h
\end{leancode}

\begin{proof}
The first is $\name{mul\_inv}$ of \texttt{PrimeTensor/Algebra.lean}.  For the
second, $1 \cdot 1^{-1} = 1$ becomes $1^{-1} = 1$ after cancelling the unit on
the left.  For the third, commute $x$ past $h$, reassociate to expose
$x x^{-1}$, cancel it, and drop the unit.
\end{proof}

The third identity is small but it is the pivot of the whole file: it says that
perturbing $x$ multiplicatively by $h$ and then comparing with $x$ returns
exactly $h$.  The response of the identity function is therefore the
perturbation itself, on the nose, with no error term --- which is
Theorem~\ref{Analysis:Derivative:thm:identity} below.

\subsection{Tangent maps}

\begin{definition}[\lean{MulTangentMap}]\label{Analysis:Derivative:def:tangent}
A \emph{tangent map} is a structure bundling a function
$\name{toFun} : \name{MulReal} \to \name{MulReal}$ with proofs that it
preserves the pivot, the multiplication and the inversion.  A $\ctor{CoeFun}$
instance lets such a structure be applied directly to an argument.
\end{definition}

\begin{leancode}
structure MulTangentMap where
  toFun : MulReal → MulReal
  map_one : toFun 1 = 1
  map_mul : ∀ a b : MulReal, toFun (a * b) = toFun a * toFun b
  map_inv : ∀ a : MulReal, toFun a⁻¹ = (toFun a)⁻¹
\end{leancode}

\begin{remark}[This is the linear map, and the axiom list is again redundant]
\label{Analysis:Derivative:rem:tangent}
In the additive theory the first-order model is required to be linear.  Here
the requirement is that the model be a morphism of the structure the carrier
actually has --- the one axiomatised as $\name{IsMulCarrier}$ in
\texttt{PrimeTensor/Structure.lean} --- and multiplicativity is what
additivity becomes.

As with that class, the axioms are not independent: $\name{map\_inv}$ follows
from the other two.  From $\name{toFun}(a^{-1})\,\name{toFun}(a) =
\name{toFun}(a^{-1}a) = \name{toFun}(1) = 1$ and uniqueness of inverses in a
carrier, $\name{toFun}(a^{-1})$ is $\name{toFun}(a)^{-1}$.  It is listed
because it is used, and listing it costs one line in each of the two instances
below.

What is \emph{not} required is any continuity or boundedness of the tangent
map.  In the additive theory over an infinite-dimensional space that
requirement is what makes the derivative well behaved; here nothing rules out a
wildly discontinuous morphism serving as a first-order model, and nothing in
this file needs to.
\end{remark}

\begin{definition}[\lean{MulTangentMap.identity} and \lean{trivial}]
\label{Analysis:Derivative:def:examples}
$\name{identity}$ is $h \mapsto h$ and $\name{trivial}$ is $h \mapsto 1$; the
$\ctor{simp}$ lemmas $\name{identity\_apply}$ and $\name{trivial\_apply}$
record their values, both by $\ctor{rfl}$.
\end{definition}

\begin{proof}[Verification]
For $\name{identity}$ all three laws are $\ctor{rfl}$.  For $\name{trivial}$,
preservation of the pivot is $\ctor{rfl}$, preservation of the product is
$1 = 1 \cdot 1$ read backwards from $\name{one\_mul}$, and preservation of the
inverse is $1 = 1^{-1}$ read backwards from
Lemma~\ref{Analysis:Derivative:lem:ratio}.
\end{proof}

\subsection{Perturbations, responses and negligibility}

\begin{definition}[\lean{ApproachesPivot} and \lean{response}]
\label{Analysis:Derivative:def:response}
A sequence $h : \name{Depth} \to \name{MulReal}$ \emph{approaches the pivot}
when it converges to $1$ in the sense of
\texttt{PrimeTensor/Analysis/Limit.lean}.  The \emph{response} of $f$ at $x$
along $h$ is the sequence
\[
  \name{response}(f,x,h)_n \defeq
    \name{ratio}\bigl(f(x \cdot h_n),\, f(x)\bigr).
\]
\end{definition}

Two features distinguish this from a difference quotient.  The perturbation
enters multiplicatively, as $x \cdot h_n$ rather than $x + h_n$, so the
statement never needs an additive structure on the domain.  And the response is
\emph{not} divided by the perturbation: it is the raw output ratio, and the
comparison against the model happens in
Definition~\ref{Analysis:Derivative:def:firstorder} through negligibility
rather than through a quotient.  This matters because dividing would require
the model to be invertible, which no tangent map is asked to be ---
$\name{trivial}$ is not.

\begin{definition}[\lean{NegligibleRelative}]
\label{Analysis:Derivative:def:negligible}
$\name{error}$ is \emph{negligible relative to} $\name{base}$ when
\[
  \forall\, g\;
  \exists\, c\;
  \forall\, n \;\text{at or after}\; c\;
  \forall\, \ell : \quad
  \name{ScaleNear}(\ell, \name{base}_n, 1)
  \;\longrightarrow\;
  \name{ScaleNear}(\name{advance}(\ell,g),\, \name{error}_n,\, 1).
\]
\end{definition}

\begin{leancode}
def NegligibleRelative (error base : MulReal.Seq) : Prop :=
  ∀ gain : Depth,
    ∃ anchor : Depth,
      ∀ n : Depth,
        Depth.AtOrAfter anchor n →
        ∀ level : Depth,
          MulReal.ScaleNear level (base n) 1 →
          MulReal.ScaleNear (Depth.advance level gain) (error n) 1
\end{leancode}

Read it as a scale shift: whatever level the base achieves at stage $n$, the
error achieves that level advanced by $g$, and $g$ may be prescribed in
advance and arbitrarily large.  Two points of quantifier order are worth
noticing.  The anchor is chosen before the level, so one anchor must serve for
every level simultaneously --- the shift is uniform in $\ell$.  And the
implication is per-stage, so at a stage where the base achieves no level at all
the condition says nothing about the error there; the definition constrains the
error only where the base is genuinely small, which for a perturbation is
eventually everywhere.

\begin{designnote}[All of this in the logarithmic chart]
\label{Analysis:Derivative:note:log}
Write $L = \log_2$, so that $L$ carries the multiplicative carrier to an
additive one, $L(1) = 0$, and $L(\name{ratio}(a,b)) = L(a) - L(b)$.  By the
computation of \texttt{PrimeTensor/Scale.lean}, closeness at level $\ell$ is
\[
  \name{ScaleWithin}(\ell,a,b)
  \iff
  \lvert L(a) - L(b)\rvert < 2^{-(\size{\ell}-1)} ,
\]
so a level is a dyadic tolerance and $\ctor{succ}$ halves it.  Under $L$:

\emph{Perturbations} approaching the pivot become sequences with
$L(h_n) \to 0$.  \emph{Responses} become
$L f(x \cdot h_n) - L f(x)$, that is, the increment of $L \circ f$ at the
argument increment $L(h_n)$.  \emph{Tangent maps} become additive maps.
And \emph{negligibility}, since $\name{advance}(\ell,g)$ adds $\size{g}$ to the
level, becomes
\[
  \lvert L(\name{base}_n)\rvert < 2^{-(\size{\ell}-1)}
  \;\Longrightarrow\;
  \lvert L(\name{error}_n)\rvert < 2^{-\size{g}} \cdot 2^{-(\size{\ell}-1)} ,
\]
which is $\lvert L(\name{error}_n)\rvert \leq 2^{-\size{g}}
\lvert L(\name{base}_n)\rvert$ up to the factor of at most two lost by rounding
the bound to the next dyadic value.  Since $g$ is arbitrary, that is exactly
$\name{error} = o(\name{base})$, with the rate quantified in powers of two
rather than by an arbitrary positive real.

Putting the four together,
Definition~\ref{Analysis:Derivative:def:hasderiv} says
\[
  L f(x \cdot h_n) - L f(x) = \tilde{D}\bigl(L(h_n)\bigr) + o\bigl(L(h_n)\bigr),
  \qquad \tilde{D} \text{ additive},
\]
which is the ordinary first-order expansion of $L \circ f \circ L^{-1}$ at
$L(x)$.  The multiplicative derivative of $f$ is the ordinary derivative of the
conjugated function; the content of the file is that this can be said without
ever leaving the multiplicative language, so that no logarithm, no real
exponent and no additive $\varepsilon$ is introduced.  The chart is a reading
aid, and nothing in this paragraph is proved in Lean.
\end{designnote}

\begin{definition}[\lean{FirstOrderEquivalent} and \lean{HasMulDerivativeAt}]
\label{Analysis:Derivative:def:firstorder}
\label{Analysis:Derivative:def:hasderiv}
Germs $\name{actual}$ and $\name{model}$ \emph{agree to first order relative to}
$\name{base}$ when $n \mapsto \name{ratio}(\name{actual}_n, \name{model}_n)$ is
negligible relative to $\name{base}$.  And $D$ \emph{is a multiplicative
derivative of $f$ at $x$} when for every perturbation $h$ approaching the
pivot, the response of $f$ at $x$ along $h$ agrees to first order, relative to
$h$, with $n \mapsto D(h_n)$.
\end{definition}

\begin{leancode}
def HasMulDerivativeAt
    (f : MulReal → MulReal) (x : MulReal) (D : MulTangentMap) : Prop :=
  ∀ h : MulReal.Seq,
    ApproachesPivot h →
    FirstOrderEquivalent h
      (response f x h)
      (fun n => D (h n))
\end{leancode}

The error between the actual response and the model is measured as their
\emph{ratio}, which is the multiplicative counterpart of their difference; the
two are the same object in the chart of
Design note~\ref{Analysis:Derivative:note:log}.

\begin{remark}[Sequential, not filtered]\label{Analysis:Derivative:rem:sequential}
The definition quantifies over sequences $h$ rather than over a neighbourhood
filter of the pivot, so it is a sequential notion --- closer to a Gateaux
derivative along all pivot-approaching sequences than to a Fréchet derivative.
Nothing here asserts that the scale shift is uniform over the choice of $h$,
only over the level for each fixed $h$.
\end{remark}

\subsection{The two computable cases}

\begin{lemma}[\lean{firstOrder\_of\_exact}]\label{Analysis:Derivative:lem:exact}
If $\name{actual}_n = \name{model}_n$ for every $n$, then they agree to first
order relative to any base.
\end{lemma}

\begin{proof}
Given a gain, take the anchor $\ctor{one}$.  For any $n$ and any level, the
ratio of the two germs at $n$ is $\name{ratio}(m,m) = 1$ by
Lemma~\ref{Analysis:Derivative:lem:ratio}, and
$\name{ScaleNear}(\ell', 1, 1)$ holds at every level $\ell'$ by reflexivity.
The hypothesis about the base is not used.
\end{proof}

\begin{theorem}[\lean{hasMulDerivativeAt\_identity}]
\label{Analysis:Derivative:thm:identity}
The identity function has derivative $\name{identity}$ at every point.
\end{theorem}

\begin{proof}
Fix $x$ and a pivot-approaching $h$.  The response at stage $n$ is
$\name{ratio}(x \cdot h_n, x)$, which is $h_n$ by
Lemma~\ref{Analysis:Derivative:lem:ratio}, and the model at stage $n$ is
$\name{identity}(h_n) = h_n$.  The two germs are equal term by term, so
Lemma~\ref{Analysis:Derivative:lem:exact} applies.
\end{proof}

\begin{theorem}[\lean{hasMulDerivativeAt\_constant}]
\label{Analysis:Derivative:thm:constant}
For every $c$, the constant function at $c$ has derivative $\name{trivial}$ at
every point.
\end{theorem}

\begin{proof}
The response at stage $n$ is $\name{ratio}(c,c) = 1$ and the model is
$\name{trivial}(h_n) = 1$; again the germs are equal term by term.
\end{proof}

\begin{remark}[Neither example exercises the machinery]
\label{Analysis:Derivative:rem:exact-only}
Both theorems go through Lemma~\ref{Analysis:Derivative:lem:exact}, whose proof
discards the base hypothesis and returns the anchor $\ctor{one}$.  So neither
computation uses $\name{advance}$, negligibility, the convergence of $h$, or
any property of the scale beyond reflexivity of $\name{ScaleNear}$: the
hypothesis $\name{ApproachesPivot}\,h$ is accepted and never read in either
proof.  These are the two cases where the model is exact, and they are here as
sanity checks on the definitions rather than as evidence that the definitions
have content.  The first genuinely approximate derivative is not computed in
this file.
\end{remark}

\begin{remark}[What is not proved]\label{Analysis:Derivative:rem:missing}
The derivative is not shown unique: nothing rules out two different tangent
maps both modelling $f$ at $x$, and the argument that would exclude it needs
the separation property still missing since
\texttt{PrimeTensor/Structure.lean}.  There is no chain rule, no product rule,
and no statement that a function with a derivative is in any sense continuous.
$\name{NegligibleRelative}$ is not shown transitive, nor closed under products,
nor related to $\name{ConvergesTo}$.  And, consistently with the header,
nothing additive is used anywhere: no subtraction, no zero, no zeroth index,
and no additive $\varepsilon$.
\end{remark}

\subsection*{Summary of the correspondence}

\begin{center}
\small
\begin{tabular}{@{}lll@{}}
\hline
Lean declaration & Kind & Here \\
\hline
\lean{Depth.advance}                   & definition & Def.~\ref{Analysis:Derivative:def:advance} \\
\lean{MulReal.ratio}                   & definition & Def.~\ref{Analysis:Derivative:def:ratio} \\
\lean{MulReal.ratio\_self}             & simp thm   & Lem.~\ref{Analysis:Derivative:lem:ratio} \\
\lean{MulReal.inv\_one}                & simp thm   & Lem.~\ref{Analysis:Derivative:lem:ratio} \\
\lean{MulReal.ratio\_mul\_base}        & theorem    & Lem.~\ref{Analysis:Derivative:lem:ratio} \\
\lean{MulTangentMap}                   & structure  & Def.~\ref{Analysis:Derivative:def:tangent} \\
\lean{CoeFun MulTangentMap}            & instance   & Def.~\ref{Analysis:Derivative:def:tangent} \\
\lean{MulTangentMap.identity}          & definition & Def.~\ref{Analysis:Derivative:def:examples} \\
\lean{MulTangentMap.trivial}           & definition & Def.~\ref{Analysis:Derivative:def:examples} \\
\lean{MulTangentMap.identity\_apply}   & simp thm   & Def.~\ref{Analysis:Derivative:def:examples} \\
\lean{MulTangentMap.trivial\_apply}    & simp thm   & Def.~\ref{Analysis:Derivative:def:examples} \\
\lean{MulDifferential.ApproachesPivot} & definition & Def.~\ref{Analysis:Derivative:def:response} \\
\lean{MulDifferential.response}        & definition & Def.~\ref{Analysis:Derivative:def:response} \\
\lean{MulDifferential.NegligibleRelative} & definition & Def.~\ref{Analysis:Derivative:def:negligible} \\
\lean{MulDifferential.FirstOrderEquivalent} & definition & Def.~\ref{Analysis:Derivative:def:firstorder} \\
\lean{MulDifferential.HasMulDerivativeAt} & definition & Def.~\ref{Analysis:Derivative:def:hasderiv} \\
\lean{MulDifferential.firstOrder\_of\_exact} & theorem & Lem.~\ref{Analysis:Derivative:lem:exact} \\
\lean{hasMulDerivativeAt\_identity}    & theorem    & Thm.~\ref{Analysis:Derivative:thm:identity} \\
\lean{hasMulDerivativeAt\_constant}    & theorem    & Thm.~\ref{Analysis:Derivative:thm:constant} \\
\hline
\end{tabular}
\end{center}

\noindent
Every declaration of \texttt{PrimeTensor/Analysis/Derivative.lean} appears
above, together with the four fields of $\name{MulTangentMap}$.  The file
contains no \lean{sorry}, no axiom, and no unproved named proposition.
Design note~\ref{Analysis:Derivative:note:log} and
Remarks~\ref{Analysis:Derivative:rem:tangent},
\ref{Analysis:Derivative:rem:sequential},
\ref{Analysis:Derivative:rem:exact-only}
and~\ref{Analysis:Derivative:rem:missing} are stated for the reader and are not
declarations of the file; in particular the logarithmic chart is a reading aid
and is not formalised.
